% Source coverage: physical PDF pp. 181--185 (printed pp. 158--162), beginning at Problems on p. 181 and ending with Reference 26 on p. 185.
% Contents: Problems 1--6; 32 displayed reference entries with identifiers 1--26 plus 3a, 4a--4d, and 8a. Figures/tables/footnotes: none.

\chapterexercisehook{18}

\chapterbackmatter{References}

\begin{enumerate}
  \setlength{\itemsep}{0.35\baselineskip}
  \item \phantomsection\label{ch18-ref-1}
  M.~Gell-Mann and F.~E.~Low, \textit{Phys. Rev.} \textbf{95}, 1300 (1954). The freedom of choosing the definition of renormalized coupling constants had been discussed a little earlier by E.~C.~G.~Stueckelberg and A.~Petermann, \textit{Helv. Phys. Acta} \textbf{26}, 499 (1953) (who also introduced the unfortunate term `renormalization group') but without an explanation of its relevance to the calculation of physical processes at very high energy or very low energy. After the work of Gell-Mann and Low renormalization group methods were further developed by N.~N.~Bogoliubov and D.~V.~Shirkov, \textit{Introduction to the Theory of Quantized Fields} (Interscience, New York, 1959): Chapter VIII, and references quoted therein. Interest in renormalization group methods in particle physics was revived in 1970 by C.~G.~Callan, \textit{Phys. Rev.} D\textbf{2}, 1541 (1970); K.~Symanzik, \textit{Commun. Math. Phys.} \textbf{18}, 227 (1970); C.~G.~Callan, S.~Coleman, and R.~Jackiw, \textit{Ann. of Phys.} (New York), \textbf{59}, 42 (1970).

  \item \phantomsection\label{ch18-ref-2}
  K.~G.~Wilson, \textit{Phys. Rev.} B\textbf{4}, 3174, 3184 (1971); \textit{Rev. Mod. Phys.} \textbf{47}, 773 (1975).

  \item \phantomsection\label{ch18-ref-3}
  J.~C.~Collins, \textit{Phys. Rev.} D\textbf{10}, 1213 (1974).

  \item[3a.] \phantomsection\label{ch18-ref-3a}
  R.~Jost and J.~M.~Luttinger, \textit{Helv. Phys. Acta} \textbf{23}, 201 (1950).

  \item \phantomsection\label{ch18-ref-4}
  This is the value quoted by the Particle Data Group in 1994. More recent calculations have been summarized by G.~Altarelli, CERN preprint CERN-TH-95/203, to be published in \textit{Proceedings of the Workshop on Physics at DA$\Phi$NE, April 1995}. These more recent values of $\alpha^{-1}(m_Z)$ range from 128.89 to 129.08.

  \item[4a.] \phantomsection\label{ch18-ref-4a}
  S.~Coleman and E.~Weinberg, \textit{Phys. Rev.} D\textbf{7}, 1888 (1973).

  \item[4b.] \phantomsection\label{ch18-ref-4b}
  L.~D.~Landau, in \textit{Niels Bohr and the Development of Physics} (Pergamon Press, New York, 1955): p.~52; and earlier works cited therein.

  \item[4c.] \phantomsection\label{ch18-ref-4c}
  S.~L.~Adler, C.~G.~Callan, D.~J.~Gross, and R.~Jackiw, \textit{Phys. Rev.} D\textbf{6}, 2982 (1972); M.~Baker and K.~Johnson, \textit{Physica} \textbf{96A}, 120 (1979).

  \item[4d.] \phantomsection\label{ch18-ref-4d}
  For a discussion and references, see J.~Glimm and A.~Jaffe, \textit{Quantum Physics -- A Functional Integral Point of View}, Second Edition (Springer-Verlag, New York, 1987): Section 21.6; R.~Fernandez, J.~Fr\"ohlich, and A.~D.~Sokal, \textit{Random Walks, Critical Phenomena, and Triviality in Quantum Field Theory} (Springer-Verlag, Berlin, 1992): Chapter 15.

  \item \phantomsection\label{ch18-ref-5}
  S.~Weinberg, in \textit{General Relativity}, eds. S.~W.~Hawking and W.~Israel, eds. (Cambridge University Press, Cambridge, 1979): p.~790.

  \item \phantomsection\label{ch18-ref-6}
  S.~Weinberg, \textit{Phys. Rev.} D\textbf{8}, 3497 (1973).

  \item \phantomsection\label{ch18-ref-7}
  The original calculation is by K.~G.~Wilson and M.~E.~Fisher, \textit{Phys. Rev. Lett.} \textbf{28}, 240 (1972); K.~G.~Wilson, \textit{Phys. Rev. Lett.} \textbf{28}, 548 (1972). For reviews, see K.~G.~Wilson and J.~Kogut, \textit{Phys. Rep.} \textbf{12C}, No.~2 (1974); M.~E.~Fisher, \textit{Rev. Mod. Phys.} \textbf{46}, 597 (1974); E.~Br\'ezin, J.~C.~Le Guillou, and J.~Zinn-Justin, in \textit{Phase Transitions and Critical Phenomena}, eds. C.~Domb and M.~S.~Green (Academic Press, London, 1975).

  \item \phantomsection\label{ch18-ref-8}
  See, e.g., P.~M.~Chaikin and T.~C.~Lubensky, \textit{Principles of Condensed Matter Physics} (Cambridge University Press, Cambridge, 1995): p.~231.

  \item[8a.] \phantomsection\label{ch18-ref-8a}
  E.~Br\'ezin, J.~C.~Le Guillou, J.~Zinn-Justin, and B.~G.~Nickel, \textit{Phys. Lett.} \textbf{44A}, 227 (1973); K.~G.~Wilson and J.~Kogut, Ref.~7, Section 8.

  \item \phantomsection\label{ch18-ref-9}
  G.~'t Hooft, \textit{Nucl. Phys.} B\textbf{61}, 455 (1973); \textit{Nucl. Phys.} B\textbf{82}, 444 (1973). The derivation presented here is a somewhat simplified version of 't Hooft's.

  \item \phantomsection\label{ch18-ref-10}
  O.~W.~Greenberg, \textit{Phys. Rev. Lett.}, \textbf{13}, 598 (1964); M.~Y.~Han and Y.~Nambu, \textit{Phys. Rev.} \textbf{139}, B1006 (1965); W.~A.~Bardeen, H.~Fritzsch, and M.~Gell-Mann, in \textit{Scale and Conformal Invariance in Hadron Physics}, ed. R.~Gatto (Wiley, New York, 1973).

  \item \phantomsection\label{ch18-ref-11}
  M.~Gell-Mann, \textit{Phys. Lett.} \textbf{8}, 214 (1964); G.~Zweig, CERN preprint TH401 (1964).

  \item \phantomsection\label{ch18-ref-12}
  D.~J.~Gross and F.~Wilczek, \textit{Phys. Rev. Lett.}, \textbf{30}, 1343 (1973).

  \item \phantomsection\label{ch18-ref-13}
  H.~D.~Politzer, \textit{Phys. Rev.}, \textbf{30}, 1346 (1973).

  \item \phantomsection\label{ch18-ref-14}
  E.~D.~Bloom \textit{et al.}, \textit{Phys. Rev. Lett.} \textbf{23}, 930 (1969); M.~Breidenbach \textit{et al.}, \textit{Phys. Rev. Lett.} \textbf{23}, 933 (1969); J.~L.~Friedman and H.~W.~Kendall, \textit{Annual Reviews of Nuclear Science} \textbf{22}, 203 (1972).

  \item \phantomsection\label{ch18-ref-15}
  A.~Zee, unpublished.

  \item \phantomsection\label{ch18-ref-16}
  G.~'t Hooft, unpublished.

  \item \phantomsection\label{ch18-ref-17}
  S.~Weinberg, \textit{Phys. Rev.} D\textbf{8}, 605 (1973).

  \item \phantomsection\label{ch18-ref-18}
  D.~J.~Gross and F.~Wilczek, \textit{Phys. Rev.} D\textbf{8}, 3633 (1973); S.~Weinberg, \textit{Phys. Rev. Lett.} \textbf{31}, 494 (1973).

  \item \phantomsection\label{ch18-ref-19}
  A similar idea had been suggested before the discovery of asymptotic freedom by H.~Fritzsch, M.~Gell-Mann, and H.~Leutwyler, \textit{Phys. Lett.} \textbf{47B}, 365 (1973).

  \item \phantomsection\label{ch18-ref-20}
  G.~Hanson \textit{et al.}, \textit{Phys. Rev. Lett.} \textbf{35}, 1609 (1975); R.~F.~Schwitters, in \textit{Proceedings of the International Conference on Lepton and Photon Interactions at High Energy at Stanford, 1975}, ed. W.~T.~Kirk (Stanford Linear Accelerator Center, Stanford, 1975): p.~5; G.~Hanson, Stanford Linear Accelerator Center Report SLAC-PUB-1814 (1976), unpublished.

  \item \phantomsection\label{ch18-ref-21}
  G.~Sterman and S.~Weinberg, \textit{Phys. Rev. Lett.} \textbf{39}, 1436 (1977).

  \item \phantomsection\label{ch18-ref-22}
  J.~Ellis, M.~K.~Gaillard, and G.~G.~Ross, \textit{Nucl. Phys.} B\textbf{111}, 253 (1976).

  \item \phantomsection\label{ch18-ref-23}
  E.~Eichten, K.~Lane, and M.~Peskin, \textit{Phys. Rev. Lett.} \textbf{50}, 811 (1983).

  \item \phantomsection\label{ch18-ref-24}
  For a review, see I.~Hinchliffe, in `Review of Particle Properties,' \textit{Phys. Rev.} D\textbf{50}, 1177 (1994): Section 25.

  \item \phantomsection\label{ch18-ref-25}
  G.~Altarelli, in \textit{Proceedings of the Rencontres de Hanoi}, CERN preprint CERN-PPE/94-71 (1994).

  \item \phantomsection\label{ch18-ref-26}
  K.~Abe \textit{et al.} (SLD collaboration), \textit{Phys. Rev.} D\textbf{51}, 962 (1995). A summary of earlier data on $Z^0$ decay into hadrons by M.~Shifman, Minnesota preprint hep-ph/9501222 (1995), gave a value $\alpha_s(m_Z)=0.125\pm0.005$, corresponding to $\Lambda\approx500$ MeV.
\end{enumerate}
